% Pursuing Stacks (cotes 134-2 to 134-8): archive cote 134-7, site page 67 (PDF page 68), typescript page 509.
% typescript. Draft compared: ps6.tex:3282-3365.
% AI-transcribed by Claude Opus 5.5 (claude-opus-5-5, adaptive thinking, effort high), Message Batches API, batch ps4 (msgbatch_012Aoc1NuskATPsyBeKSYK7J), 2026-11-12; not yet checked by eye.
\scanpage{134-7/67}
\mspage{509}
K(\hand{$\pi$},n). As any object in H\over{o}{\unsure{$\theta$}}tab \xout{\unsure{can} be written} is isomorphic to a
\margin{direct argument}product of objects $\hand{\pi}_n[n]$, it follows that \inserted[typed]{in order} \xout{\illegible{1}} to check the adjunction
formula between LH. and E it is enough to do so for objects in Hotab
of the type $\hand{\pi}[n]$, which is the ``familiar fact'' just recalled. The nota-%
tion $K_{\hand{\pi}}$ in (2) is meant to suggest the Eilenberg-Maclane K(\hand{$\pi$},n) objec\unsure{t}
generalized by the objects $K_{\hand{\pi}}(L.)$ , and also to recall that we recover
the homology invariants of \xout{K(L.)} L. from $K_{\hand{\pi}}(L.)$ \inserted{via the $\pi_i$ invariants,} by the formula
\[
(4) \qquad \hand{\pi}_i(K_{\hand{\pi}}(L.)) \hand{\simeq} H_i(L.) \quad ,
\]
which implies by the way that the functor $K_{\hand{\pi}}$ is ``conservative'', i.e.
a map in Hotab which is transformed into an isomorphism is an \unsure{isomorphism.}
This should not be confused with the stronger property of being fully
faithful, or equivalently of the left adjoint LH. being a lo\over{\unsure{d}}{c}alization
functor, or equivalently still, the adjunction morphism
\[
(5) \qquad LH.(K_{\hand{\pi}}(L.)) \hand{\to} L.
\]
being an isomorphism in H\over{o}{\u{o}}tab, which is definitely false !

The other adjunction morphism
\[
(6) \qquad X \hand{\to} K_{\hand{\pi}}(LH.(X)) \overset{\mathrm{def}}{=} X_{ab}
\]
is still more interesting, it's effect on the homotopy invariants $\hand{\pi}_i$ \deleted{are}\inserted{are}
the Hurewicz \xout{maps} homomorphisms
\[
(7) \qquad \hand{\pi}_i(X) \hand{\to} H_i(X) \overset{\mathrm{def}}{=} \hand{\pi}_i(X_{ab}) \deleted{\overset{\mathrm{def}}{=}}\hand{\simeq} H_i(LH.(X)) \quad ;
\]
\deleted{\unsure{thus}} introducing the homotopy fiber of (6) (in the case of pointed
homotopy types) and denoting by $\hand{\gamma}_i(X)$ it's homotopy invariants, we get
the exact sequence of J.H.C. Whitehead (as recalled in a letter from R.
Brown I just got)
\[
\hand{(7')} \qquad \dots \hand{\to} \hand{\gamma}_i(X) \hand{\to} \hand{\pi}_i(X) \hand{\to} \hand{\widetilde{\phantom{H}}}\!\!H_i(X) \hand{\to} \hand{\gamma}_{i-1}(X) \hand{\to} \quad \dots \quad ,
\]
where $\hand{\widetilde{\phantom{H}}}\!\!H_i$ denotes the ``reduced'' homology group of a \uline{pointed} homotopy
type, equal to $H_i$ for $i\hand{\neq}0$ and to Coker $(H_o(pt) \hand{\to} H_o(X)) \hand{\simeq} H_o(X)/\underline{Z}$ for
i = 0.

The case of pointed homotopy types seems of importance for schema-%
tic homotopy types, and deserves some extra mention and care. We may
factor diagram (2) into
\margin{$\widetilde{DP}$ (circled)}
\diagram{
\begin{tikzcd}
(8) \qquad \hand{\Delta^{\wedge}} \ar[r, shift left, "\hand{\alpha}"] & \hand{\Delta^{\wedge\bullet}} \ar[l, shift left, "\hand{\beta}"] \ar[r, shift left, "\hand{\widetilde{W}}"] & \hand{\Delta^{\wedge}}_{ab} \ar[l, shift left, "\hand{\widetilde{DP}}"]
\end{tikzcd} \quad ,
}
where $\hand{\Delta^{\wedge\bullet}}$ is the category of \uline{pointed} semisimplicial complexes, \hand{$\beta$} the
\xout{\illegible{1}} forgetful functor from these to nonpointed complexes, and \hand{$\alpha$} it's
left adjoint, which may be interpreted as

\margin{67}
